\subsection{Topic 6: Concurrency \& Parallelism – Exercise Sheet [Solutions]}
\subsection*{Q\#1: A supermarket has sensors on each entrance that detect when a customer enters or leaves. These
sensors send a signal to a central computer whenever a customer enters or leaves. This computer,
amongst other things, keeps a running count of how many people are in the shop at any one time. This
is used to control a traffic light system that only allows customers to enter the store when the number of
customers is below a safe limit. Last week, the store was empty - but the system refused to let anyone
in. What went wrong? How would you fix it?} 
\textcolor{blue}{
The sensors at the entrances send signals to the central computer system when a customer enters or
leaves the store. On receiving a signal, the computer either increments or decrements (depending on the
type of signal) the count of how many customers are inside the supermarket. Now, this increment /
decrement operation, as it is given in the question, has not been implemented properly i.e. the counter
is a shared resource and the increment / decrement operations do not take a lock before updating it,
which means that we can have race conditions. It is quite possible that the value of the counter will be
incorrect after a number of these operations, causing the system to malfunction i.e. the store is empty
and the system thinks that there are customer inside the shop.
The obvious solution will be to implement a locking mechanism e.g. mutual exclusion before updating
the value of counter. 
}
\subsection*{Q\#2: An e-commerce application sells shoes. The application stores the current stock level of each
shoe in memory. When a pair of shoes is purchased a thread is created. The thread i) checks the current
stock level for that shoe ii) calculates the new value and iii) updates the stock level. The stock level
recorded is usually correct but occasionally the application shows an incorrect value.}
\begin{itemize}
    \item{a)} Explain the (technical) reason why the application occasionally shows an incorrect value.
    \item{b)} Explain what changes should be made to the program to correct this behaviour.
    \item{c)} After these changes, the system works correctly. However, the performance has degraded.
\end{itemize}
\textcolor{blue}{
Explain clearly why the performance will have got worse.
The current stock level for each shoe is stored in memory and is a shared resource. Whenever a pair of
shoes is purchased, a thread is created, and there is a race condition where 2 threads could be updating
the same shoe stock level value at the same time. 
\newline
In order to correct this problem, we need to identify the critical section and ensure that only one thread
can be in this section at a given time. We can make use of locking mechanisms, such as a mutex to
implement this solution. 
\newline
Once the solution has been implemented, performance will most obviously be affected because of a)
the extra cost of protecting the critical section and b) because operations now need to be sequential.
}
\subsection*{Q\#3: Consider the following program:}
% \documentclass{article}
% \usepackage{listings}
% \usepackage{xcolor}

\lstset{
  language=Python,
  basicstyle=\ttfamily\small,
  keywordstyle=\color{blue},
  commentstyle=\color{gray},
  stringstyle=\color{red},
  showstringspaces=false,
  breaklines=true,
  frame=single,
}

\begin{lstlisting}
boolean blocked [2];
int turn;
void P (int id)
{
    while (true) {
        blocked[id] = true;
        while (turn != id) {
            while (blocked[1-id])
                /* do nothing */;
            turn = id;
        }
        /* critical section */
        blocked[id] = false;
        /* remainder */
    }
}
void main()
{
    blocked[0] = false;
    blocked[1] = false;
    turn = 0;
    parbegin (P(0), P(1));
}
\end{lstlisting}

The above software solution to the mutual exclusion problem for two processes has been suggested in
the literature. Do you agree that the proposed solution is correct? If yes, explain with the help of
supporting arguments and examples. If no, find a counter-example that demonstrates that this solution
is incorrect and does not preserve mutual exclusion.
\begin{figure}[htp]
  \centering
  \includegraphics[width=20cm]{unit-22/figures/01.png}
  \caption*{currency}
\end{figure}

\subsection*{Q\#4: Your friend Jane needs to run a simulation for her thesis and her adviser wants her to run it for a
fixed problem size. Jane can make 90\% of the program parallel, with 10\% of it being sequential.}
\begin{itemize}
    \item{a)} What is the maximum speedup Jane can expect on 10 processors? 
    \newline
    Speedup = \( \frac{1}{(1 - f) + \frac{f}{N}} \) = \( \frac{1}{(1 - 0.9) + \frac{0.9}{10}} \) = \( \frac{1}{0.1 + 0.09} = 5.26x \)
    \item{b)} What would be the maximum speedup on an infinite number of processors?
    \newline
    Max Speedup = \( \frac{1}{(1 - 0.9) + \frac{0.9}{\infty}} = \frac{1}{0.1 + 0} = 10x \)
\end{itemize}
\subsection*{Q\#5: You asked your friend Mike to confirm the results that you obtained from running your code,
having 80\% of it in parallel and 20\% sequential. Your friend confirmed the same result but got it 2.5x
times faster.}

a) How many processors did he use?
\( 2.5 = \frac{1}{(1 - 0.8) + \frac{0.8}{N}} \)
Solving gives: \( N = 4 \)

b) Will the speed-up be doubled, if we double the number of processors used?
Speedup = \( \frac{1}{0.2 + \frac{0.8}{8}} = 3.33x \), No

\subsection*{Q\#6: A program speeds up by 4x when five processors are used. Determine the percentage of the code
running in parallel.}  

Speedup = \( \frac{1}{(1 - f) + \frac{f}{5}} = 4 \)

Solving gives: \( f = 0.9375 \Rightarrow 93.75\% \) of the code is parallel.

\subsection*{Q\#7: Assume a system where a particular task can be parallelized to 96\%. Calculate the speedup
gained by using 16 processors instead of one, using Amdahl’s law. If the sequential version of the task
takes 140 seconds, how much time the parallel version will take on 16 processors?}

Speedup = \( \frac{1}{(1 - 0.96) + \frac{0.96}{16}} = \frac{1}{0.04 + 0.06} = 10x \)
Sequential version time: 140 seconds  \\
Parallel Time = Sequential Time / Speedup = \( \frac{140}{10} = 14 \) seconds

